\documentclass[10pt,a4paper]{amsart}
\usepackage[margin=19mm]{geometry}
\usepackage{amsmath,amssymb,mathtools}
\usepackage[scr=rsfs,cal=euler]{mathalfa}
\usepackage{xfrac}
\usepackage[unicode,hidelinks,bookmarks=false]{hyperref}
\newcommand{\ie}{i.e.\ }
\newcommand{\cf}{cf.\ }
\newcommand{\resp}{respectively\ }
\newcommand{\Subsection}{\subsection}
\providecommand{\nobreakdash}{\nobreak\hbox{-}}
\setlength{\emergencystretch}{2em}
\multlinegap0pt
\usepackage{amssymb}
\usepackage{tikz}
\newtheorem{corollary}{Corollary}
\newtheorem{proposition}{Proposition}
\newtheorem{lemma}{Lemma}
\newcommand{\HH}{\mathcal H}                 % Hankel transform
\newcommand{\R}{\mathbb R}
\newcommand{\T}{\mathbb T}
\newcommand{\Z}{\mathbb Z}
\DeclareMathOperator*{\slim}{s\text{-}lim}
\title{$L^p$ bounds for wave operators with critical electromagnetic potentials}
\author{Piero D'Ancona, Xitao Gao, Junyong Zhang}
\date{Source: arXiv:2609.07464v1 (CC BY 4.0)}
\begin{document}
\maketitle

\noindent\emph{Source and licence.} $L^p$ bounds for wave operators with critical electromagnetic potentials. Version arXiv:2609.07464v1 (CC BY 4.0), distributed by arXiv under the Creative Commons Attribution 4.0 International licence, http://creativecommons.org/licenses/by/4.0/, as recorded in the arXiv licence metadata for this exact version and shown on the abstract page https://arxiv.org/abs/2609.07464v1. Full licence notice: \texttt{licenses/arxiv-2609.07464-CC-BY-4.0.txt}.\par\smallskip

\noindent\textit{Editorial scope.} Counted unit: the article's proposition prop:em-wave (source0.tex lines 1309--1323). The article's complete original proof of it is delivered with it (source0.tex lines 1508--1613, one \texttt{proof} environment of 106 lines, which the article places away from the statement). No mathematics is written, repaired or re-derived: the statement and the proof are the article's own lines, in the article's own notation, and no retained line is substituted or re-flowed.  Retained uncounted as the source-local context the proof needs: lines 302--307; lines 738--741; lines 747--756; lines 750--754; lines 1334--1345; lines 1398--1411; lines 1403--1408; lines 1479--1500; lines 1489--1492.  Nothing was omitted from any retained span, so the package carries no omission record.  No retained line contains a citation, so the wrapper carries no bibliography and no bibliography entry.  No retained line contains a figure environment, an image-inclusion command or drawing code, so no image is delivered and the package declares no asset dependency.  Editorial formatting only, in addition: an AMS \texttt{amsart} preamble that restates the article's own declarations and macros used by the retained lines, namely the environments \texttt{corollary}, \texttt{proposition}, \texttt{lemma} on separate counters, together with the standard packages those lines need.  The retained environments print in this document as Corollary 1, Proposition 1, Lemma 1 in order of appearance, whereas the article prints the same items with the numbering of its own full document.  The complete native source of the article as distributed in the arXiv e-print for this exact version is bundled unchanged as \texttt{sources/209630/source0.tex}.
\begin{corollary}[Friedrichs $L^{p}$ bounds]
\label{cor:AB-unweighted}
  For every $\alpha\in\R$ and $1<p<\infty$, the operators
  $W_\pm,W_\pm^*$ are bounded on $L^p(\R^2)$, with constants
  locally bounded in $\alpha$. 
\end{corollary}

\begin{equation}\label{eq:em-K}
  \mathcal K_B\left(\sum_n f_n(r)\psi_n(\theta)\right)
  =\sum_n(\HH_{s_n}f_n)(r)\psi_n(\theta).
\end{equation}

\begin{proposition}\label{prop:em-L2}
  For two such angular operators $B_0,B_1$, the ordinary
  wave operators exist and are unitary. More precisely,
  \begin{equation}\label{eq:em-L2}
    W_\pm(H_{B_1},H_{B_0})
    =\mathcal K_{B_1}e^{\pm i\pi\sqrt{B_1}/2}
    e^{\mp i\pi\sqrt{B_0}/2}\mathcal K_{B_0}.
  \end{equation}
  No relation between the two angular eigenbases is needed.
\end{proposition}

  \begin{equation}
    W_\pm(H_{B_1},H_{B_0})
    =\mathcal K_{B_1}e^{\pm i\pi\sqrt{B_1}/2}
    e^{\mp i\pi\sqrt{B_0}/2}\mathcal K_{B_0}.
  \end{equation}

\begin{lemma}\label{lem:em-angular}
  Let $e_k=(2\pi)^{-1/2}e^{ik\theta}$, and suppose that a
  unitary angular operator $J$ satisfies
  $Je_k=e_k(1+r_k)$, where
  \begin{equation}\label{eq:em-summability}
    \sum_{k\in\Z}\|r_k\|_{H^3(\T)}^2<\infty.
  \end{equation}
  Then $J,J^*$ and
  $\mathcal F_\pm^{-1}J\mathcal F_\pm$,
  $\mathcal F_\pm^{-1}J^*\mathcal F_\pm$ are bounded on
  $L^p(\R^2)$ for every $1<p<\infty$.
\end{lemma}

\begin{lemma}\label{lem:em-eigenbasis}
  Under the assumptions of Proposition~\ref{prop:em-wave}, there
  is an orthonormal eigenbasis $\{\phi_k\}_{k\in\Z}$ of
  $B_{\alpha,a}$ such that, writing
  $\phi_k=e_k b_k$ and $B_{\alpha,a}\phi_k=\lambda_k\phi_k$,
  \begin{equation}\label{eq:em-basis-estimate}
    |\lambda_k-(k+\alpha)^2|\le\|a\|_\infty,
    \qquad
    \|b_k-1\|_{H^3(\T)}
    \le\frac{C_{\alpha,a}}{1+|k|}.
  \end{equation}
  Consequently, $Je_k=\phi_k$ satisfies
  Lemma~\ref{lem:em-angular}.
\end{lemma}

  \begin{equation}
    |\lambda_k-(k+\alpha)^2|\le\|a\|_\infty,
    \qquad
    \|b_k-1\|_{H^3(\T)}
    \le\frac{C_{\alpha,a}}{1+|k|}.
  \end{equation}

\begin{proposition}[Spectral transplantation]
\label{prop:em-transplantation}
  Let $a$ and $\alpha$ be as in
  Proposition~\ref{prop:em-wave}. Choose the eigenbasis of
  Lemma~\ref{lem:em-eigenbasis} and set
  \begin{equation*}
    \sigma_k=|k+\alpha|,\qquad \tau_k=\sqrt{\lambda_k},
    \qquad Je_k=\phi_k.
  \end{equation*}
  The operator
  \begin{equation}\label{eq:em-S}
    S\left(\sum_k f_ke_k\right)
    =\sum_k(\HH_{\tau_k}\HH_{\sigma_k}f_k)\phi_k
  \end{equation}
  is unitary on $L^2(\R^2)$. Both $S$ and $S^*$ extend
  boundedly to every $L^p(\R^2)$, $1<p<\infty$, and are
  mutually inverse there. For every bounded Borel function
  $m$ on $[0,\infty)$,
  \begin{equation*}
    m(H_{\alpha,a})S=Sm(H_\alpha)\qquad\text{on }L^2.
  \end{equation*}
\end{proposition}

  \begin{equation}
    S\left(\sum_k f_ke_k\right)
    =\sum_k(\HH_{\tau_k}\HH_{\sigma_k}f_k)\phi_k
  \end{equation}

\begin{proposition}[Constant magnetic coefficient]
\label{prop:em-wave}
  Let $\alpha\in\R\setminus\frac12\Z$ and let
  $a\in W^{3,\infty}(\T;\R)$ satisfy $B_{\alpha,a}\ge0$.
  The wave operators
  \begin{equation}\label{eq:em-wave}
    W_\pm^{\alpha,a}
    =\slim_{t\to\pm\infty}e^{itH_{\alpha,a}}e^{it\Delta}
  \end{equation}
  exist and are unitary on $L^2(\R^2)$. For every
  $1<p<\infty$, both $W_\pm^{\alpha,a}$ and their
  adjoints extend boundedly to $L^p(\R^2)$ and are mutually
  inverse there. The same assertions hold for the wave
  operators of the pair $(H_{\alpha,a},H_\alpha)$.
\end{proposition}

\begin{proof}[Proof of Proposition~\ref{prop:em-wave}]
  Use $\sigma_k,\tau_k,J,S$ from
  Proposition~\ref{prop:em-transplantation}.
  To simplify notation within this proof,
  following the general construction~\eqref{eq:em-K}, write
  \begin{equation*}
    \mathcal K_a=\mathcal K_{B_{\alpha,a}},\qquad
    \mathcal K_\alpha=\mathcal K_{B_{\alpha,0}},\qquad
    \mathcal K_0=\mathcal K_{D_\theta^2}.
  \end{equation*}
  More explicitly, if $\phi_k$ is the eigenbasis from
  Lemma~\ref{lem:em-eigenbasis}, then
  \begin{equation*}
    \begin{aligned}
      \mathcal K_a\Bigl(\sum_k f_k\phi_k\Bigr)
        =\sum_k(\HH_{\tau_k}f_k)\phi_k,\quad & \quad
      \mathcal K_\alpha\Bigl(\sum_k f_ke_k\Bigr)
        =\sum_k(\HH_{\sigma_k}f_k)e_k,\\
      \mathcal K_0\Bigl(\sum_k f_ke_k\Bigr)
        &=\sum_k(\HH_{|k|}f_k)e_k.
    \end{aligned}
  \end{equation*}
  Here the subscript $0$ denotes the free angular operator
  $D_\theta^2$, in parallel with $H_0=-\Delta$.
  All identities between these operators are initially on
  $L^2$. Formula~\eqref{eq:em-S} gives
  \begin{equation*}
    S=\mathcal K_aJ\mathcal K_\alpha.
  \end{equation*}
  Proposition~\ref{prop:em-transplantation} supplies the
  $L^p$ bounds for $S$ and $S^*$.

  Next put $V_\alpha=\mathcal K_\alpha\mathcal K_0$.
  Corollary~\ref{cor:AB-unweighted} implies that
  $V_\alpha,V_\alpha^*$ are bounded on every $L^p$ in
  question. Indeed, the partial wave formula differs from
  $V_\alpha$ by the angular phases
  $e^{\pm i\pi(\sigma_k-|k|)/2}$. Each phase sequence is
  constant on each sufficiently large positive or negative
  tail. The corresponding multiplier is a linear combination
  of the angular Riesz projections and finitely many angular
  projections, all bounded on $L^p$. It follows that
  \begin{equation}\label{eq:em-U}
    U=SV_\alpha=\mathcal K_aJ\mathcal K_0
  \end{equation}
  and $U^{-1}=U^*$ are bounded on $L^p$.

  Define the diagonal angular multipliers
  \begin{equation}\label{eq:em-D}
    D_\pm e_k
    =e^{\pm i\pi(\tau_k-|k|)/2}e_k.
  \end{equation}
  By \eqref{eq:em-basis-estimate},
  $\tau_k-\sigma_k=O(|k|^{-1})$. Consequently, the sequence
  in \eqref{eq:em-D} equals
  $e^{\pm i\pi\alpha/2}$ on the positive tail and
  $e^{\mp i\pi\alpha/2}$ on the negative tail, up to an
  $\ell^2(\Z)$ remainder. The constant tails give angular
  Riesz projections, and the remainder again has an $L^1$
  convolution kernel. This proves boundedness of $D_\pm$
  and $D_\pm^{-1}$ on $L^p$.

  The Fourier transform in polar coordinates can be written
  for $f=\sum_k f_ke_k$
  \begin{equation}\label{eq:em-free-Fourier}
    \mathcal F_\pm\Bigl(\sum_k f_ke_k\Bigr)
        =\sum_k e^{\mp i\pi|k|/2}
          (\HH_{|k|}f_k)e_k =
      \mathcal F_\pm
        =e^{\mp i\pi|D_\theta|/2}\mathcal K_0.
  \end{equation}
  Proposition~\ref{prop:em-L2}, applied with $B_0=D_\theta^2$,
  now yields the exact factorization
  \begin{equation}\label{eq:em-factorization}
    W_\pm^{\alpha,a}
    =U D_\pm\mathcal F_\pm^{-1}J^*\mathcal F_\pm.
  \end{equation}
  For clarity, define $T_\pm e_k=e^{\pm i\pi\tau_k/2}e_k$.
  The diagonal operators commute with $\mathcal K_0$, and
  \eqref{eq:em-free-Fourier} gives
  \begin{equation*}
    U D_\pm\mathcal F_\pm^{-1}J^*\mathcal F_\pm
      =\mathcal K_aJ T_\pm J^*\mathcal F_\pm =
      \mathcal K_a e^{\pm i\pi\sqrt{B_{\alpha,a}}/2}
        \mathcal F_\pm.
  \end{equation*}
  The last expression is exactly \eqref{eq:em-L2} for the
  free comparison operator. This also checks both signs of
  the scattering phases.

  Lemmas~\ref{lem:em-angular} and \ref{lem:em-eigenbasis}
  control the Fourier conjugate of $J^*$ and its inverse.
  Every factor in \eqref{eq:em-factorization} is therefore
  bounded and invertible on $L^p$. Proposition~\ref{prop:em-L2}
  gives existence and unitarity on $L^2$, and density extends
  the inverse identities to $L^p$. Finally, the same
  proposition, or the chain rule for these unitary wave
  operators, gives
  \begin{equation*}
    W_\pm(H_{\alpha,a},H_\alpha)
    =W_\pm^{\alpha,a}(W_\pm^\alpha)^*,
    \qquad W_\pm^\alpha=W_\pm(H_\alpha,-\Delta).
  \end{equation*}
  Corollary~\ref{cor:AB-unweighted} proves the asserted
  bounds for this pair as well.
\end{proof}

\end{document}
